% Transcription — fonds Alexandre Grothendieck, shelfmark 150, batch 2 (pages 21-40).
% Established from the facsimile archives/batches/150/batch-02.pdf (a copy of
% 150.pdf fetched through the project's relay,
% https://grothendieck-relay.onrender.com/source/150.pdf; PDF page k =
% archivists' page 20+k, no cover sheet), per the skill transcribe-grothendieck.
% The folder is 98 pages in five batches (1-20, 21-40, 41-60, 61-80, 81-98);
% this is the second, done in parallel with the others, aligned afterwards on
% batch 1 (transcripts/150/batch-01.fr.tex).
%
% Pass: Opus 5.5 (claude-opus-5-5), 2026-09-26 — first pass, unchecked against the pages by a human.
%
% Pages skipped: none. Every page carries his mathematics. Page 37 is written
% only on its upper third (the rest blank).
%
% Pages summarised: none. Two cancelled blocks are given in a \note rather
% than transcribed line by line: the foot of page 23 (barred by long
% diagonals and rewritten at the head of page 24) and the foot of page 31 (a
% first writing of system (18), barred and rewritten on page 32).
%
% His pagination and numbered runs. The batch continues the run of batch 1
% headed « Espaces stratifiés (étude par voisinages coniques) », sheets 1)-5)
% on pp. 11-19: here the sheets are numbered top left 6 (p. 21), 7 (23),
% 8 (25), 9 (27), 10 (30), 11 (31), 12 (32), 13 (34), 14 (36), 15 (38),
% 16 (40); the unnumbered pages 22, 24, 26, 28, 29, 33, 35, 37, 39 are their
% continuations (probably versos). Page 22 carries an overwritten, illegible
% figure top left. Page 29 (paragraph 7, unnumbered) is continued by
% sheet 11 (p. 31), while sheet 10 (p. 30) continues page 28: the leaves were
% scanned in the order 29, 30. Paragraphs: circled (5) « Germes de voisinages
% côniques, ou cylindriques. Transversalité » (p. 21, taking up the cancelled
% (5) of batch 1's p. 20), circled (6.) « Relations avec stratifications de
% type I » (p. 27), 7. « Stratifications et voisinages côniques » (p. 29),
% 8. « Structure cônique sur un espace stratifié » (p. 34), 9. « Description
% d'une stratification cônique de type I » (p. 37); p. 38 opens, without a
% number, « Stratification équisingulière de type I : présentation standard »
% with conditions (1), (2), (3). Numbered formulas: (14) a)-b), (14'), (15),
% (16), (16'), (17)-(23), and conditions (e), (f); (20) has circled (a)-(d).
% His (14) and (15) repeat numbers already used in batch 1 (pp. 18, 20) for
% other formulas; they are kept as written, with a \note. No date on these
% pages.
%
% What the batch is about. Pages 21-26: germs of conical and cylindrical
% neighbourhoods (T ~ T' when lambda T = lambda' T'); compatibility of a
% closed part with a germ; equisingular germs; the equivalence of categories
% between systems (X~*, dot X, E, Y, f : dot X -> Y) and (X, Y, T) with
% X = X~* ⊔_{dot X} Y; the corollary describing compatible closed parts
% (and « cônique strict » ones, Y' = Im(dot X' -> Y)); compatibility of a
% cylindrical germ E around a bordant X_0 with a conical germ T (a)-c),
% condensed into a « Condition de compatibilité » : T∩E ≃ T_0 × I carrying the
% product conical structure); a Proposition describing such E in terms of
% (X~_0, E_X~0, Y_0, E_Y0); transversality of a finite family of cylindrical
% neighbourhoods (E_J ≃ B_J × I^J) and of such a family with a conical T.
% Pages 27-28, 30: precylindrical stratifications of type I (bordant
% X_j' -> X_j for j' a predecessor, transversality, the interval [i, j]
% anti-isomorphic to P(S): local structure « cubique », a corner of a cube),
% cylindrical stratifications, and the added axiom X_i ∩ X_j = ∪_{k<=i,j} X_k.
% Pages 29, 31-35: the recursion — for I_0 the minimal elements, a strict
% conical neighbourhood T_{I_0} of X_{I_0} transverse to the (B_alpha, E_alpha)
% and to the X_i; a theorem: systems (17) (X, (X_i), (B_alpha, E_alpha),
% (T_beta)) are equivalent to systems (18) (X*, dot X_beta, B*_alpha, X*_i;
% X_beta, B_beta,alpha; dot X_beta -> X_beta) under conditions (20), (e), (f),
% via the reconstruction (19) by mapping cones C(p_beta); then the definition
% by recursion on card I of a conical structure on a stratified space
% ((a), (22), (23), (b)). Page 36-37: remarks (existence under
% non-singularity, independence of the filtration of I, replacing the B_alpha
% by a cylindrical stratification indexed by P(Lambda)°). Pages 38-40: the
% « présentation standard » of an equisingular stratification of type I as a
% functor Sigma : Drap(I) -> M (a model topological category) — fibrant
% face maps for initial segments, boundary embeddings for the others,
% cartesian squares Sigma_{d ∪ d'} over Sigma_i; reduction to the data
% Sigma_i, Sigma_ij, Sigma_ijk = Sigma_ij ×_{Sigma_j} Sigma_jk -> Sigma_ik,
% and their repackaging as Sigma_* (graded by I), Sigma_** with source s
% (fibration), target b (boundary map) and identity i.
%
% Notation fixed once for the batch, following batch 1: \dot T, \dot X for
% boundaries, \overset{\circ}{T} for interiors, \mathbf{I} for [0,1] (his
% heavy I, used both as the interval and as the monoid acting, « I. X »),
% \mathcal C(p) for the mapping cone, X^{*}, X_i^{*}; \tilde X, \tilde X^{*}
% as written; \mathcal E (his curly E) for cylindrical germs and E for their
% representatives; \mathcal T for a conical germ; \Sigma_d for d a flag,
% \Sigma_{*}, \Sigma_{**}, \underline b, \underline s, \underline i as he
% underlines them; \mathfrak P(S) for the power set. Spelling: he writes
% « cônique » with a circumflex throughout these pages (it is plain on
% pp. 21, 24, 29, 34), and it is kept so, although batch 1 normalised it to
% « conique »; « côniquement » likewise. His underlinings in running text are
% \emph{}; circled numerals and letters are given as (n) with a \note.
%
% Hardest pages: 27 (conditions b)-c) of a precylindrical stratification,
% overwritten and hemmed in by sideways marginalia), 22 and 29 (sideways
% margin notes), 38-39 (the long crosswise notes in both margins of the
% « présentation standard »). The formulas and numbered statements carry;
% much of the sideways marginalia does not and is left \ill{}.

\documentclass[11pt,a4paper]{article}
% Le préambule commun à toutes les transcriptions du fonds Grothendieck.
%
% Il tient en une idée : **l'incertitude se compose, elle ne se résout pas.**
% Une transcription qui lisse un mot douteux a détruit la seule chose pour
% laquelle on la faisait — un lecteur doit pouvoir savoir, à chaque endroit,
% ce qui était sur le feuillet et ce qui a été ajouté. Les six macros qui
% suivent sont donc l'appareil critique, et non de la décoration.
%
% Les mêmes commandes sont reconnues par `scripts/render.mjs`, qui produit la
% vue de lecture du site. Une macro ajoutée ici doit l'être aussi là-bas,
% faute de quoi le rendu échouera bruyamment — ce qui est voulu : un rendu
% silencieusement faux est pire qu'un rendu absent.

\ProvidesPackage{grothendieck}[2026/08/08 Transcriptions du fonds Grothendieck]

\usepackage{amsmath,amssymb,amsthm}
\usepackage{mathrsfs}
\usepackage{stmaryrd}
% Les diagrammes commutatifs sont partout dans ce fonds : c'est la langue
% même de la géométrie algébrique de Grothendieck, et une transcription qui
% les rendrait en prose ne transcrirait rien.
\usepackage{tikz-cd}
% Français par défaut — c'est la langue des feuillets — mais l'anglais est
% chargé aussi : la traduction et le résumé sont des documents anglais, et la
% typographie française (espaces avant les deux-points) y serait une faute.
% Ces documents basculent par \selectlanguage{english} en tête de corps.
\usepackage[english,french]{babel}
\usepackage{fontspec}
\usepackage{geometry}
\geometry{a4paper,margin=25mm,marginparwidth=28mm}
\usepackage{marginnote}
\usepackage{xcolor}
% ulem, not soul. `soul`'s \st and \ul are built on a character-by-character
% scan that XeLaTeX's font machinery breaks: the document fails with an
% undefined control sequence rather than degrading. ulem does the same two jobs
% and is Unicode-engine safe. `normalem` stops it hijacking \emph, which the
% transcriptions use for Grothendieck's own emphasis.
\usepackage[normalem]{ulem}
\usepackage{hyperref}
\hypersetup{colorlinks=true,linkcolor=black,urlcolor=black,pdfborder={0 0 0}}

% --- Métadonnées du lot -----------------------------------------------------
% Reprises telles quelles par le script de rendu, qui les affiche en tête de
% la vue de lecture. Elles fixent ce que le fichier prétend couvrir : sans
% elles, une transcription flotte sans point d'attache dans un fonds de seize
% mille feuillets.

\newcommand{\folder}[1]{\gdef\grothFolder{#1}}
\newcommand{\batch}[1]{\gdef\grothBatch{#1}}
\newcommand{\leaves}[2]{\gdef\grothFirst{#1}\gdef\grothLast{#2}}
\newcommand{\foldertitle}[1]{\gdef\grothFoldertitle{#1}}
\gdef\grothFolder{?}\gdef\grothBatch{?}\gdef\grothFirst{?}\gdef\grothLast{?}\gdef\grothFoldertitle{}

% --- Le résumé --------------------------------------------------------------
%
% Il ouvre la lecture modernisée, sous le titre « Résumé », et il est composé
% en petit. Ce n'est pas un ornement : c'est le seul passage du document qui
% ne soit pas des mathématiques, et un lecteur doit pouvoir le reconnaître
% avant de l'avoir lu — puis le sauter s'il connaît déjà le sujet, ou s'y
% arrêter s'il ne le connaît pas. Le filet qui le suit marque la reprise du
% corps.

\newenvironment{resume}
  {\small\setlength{\parskip}{0.45em}}
  {\par\medskip\hrule\medskip}

% --- Mots-clés ---------------------------------------------------------------
%
% En anglais, en clôture du résumé de la lecture modernisée : le vocabulaire
% moderne sous lequel chercher ce que les feuillets construisent. Le manifeste
% du site (`npm run manifest`) les extrait pour étiqueter la cote entière —
% cette ligne est la source unique des tags, il n'y a pas de fichier à côté.
\newcommand{\keywords}[1]{\par\smallskip\noindent
  {\footnotesize\textsc{Keywords} --- \textit{#1}}\par}

% --- L'appareil critique ----------------------------------------------------

% Le numéro de feuillet, en marge. C'est l'ancre qui permet de régler un
% désaccord sur une formule en regardant *un* feuillet plutôt qu'une chemise
% de six cents. Le site s'en sert aussi pour tourner les pages du fac-similé
% quand on fait défiler la transcription.
\newcommand{\leaf}[1]{%
  \marginnote{\footnotesize\color{gray!70}\textsf{#1}}%
  \label{leaf:#1}\ignorespaces}

% Illisible : on ne devine pas. Un mot inventé qui se lit comme les autres est
% le pire résultat possible.
\newcommand{\ill}{\textcolor{red!60!black}{[\,\ldots\,]}}

% Lecture douteuse : le mot est donné, et signalé comme incertain.
\newcommand{\uncertain}[1]{\uline{#1}}

% Ajout de l'éditeur — une lettre manquante, un mot sous-entendu.
\newcommand{\add}[1]{\textcolor{blue!50!black}{[#1]}}

% Biffé par Grothendieck lui-même. Ce qu'il a rayé fait partie du document :
% une rature dit souvent où la pensée a changé de direction.
\newcommand{\struck}[1]{\sout{#1}}

% Note du transcripteur, sur l'état matériel du feuillet.
\newcommand{\note}[1]{\par\smallskip{\small\color{gray!60!black}\textit{[#1]}}\par\smallskip}

% Note marginale de Grothendieck, distincte du corps du texte.
\newcommand{\marginal}[1]{\par\smallskip\noindent\hspace{1em}%
  {\small\color{gray!60!black}#1}\par\smallskip}

% --- Titre ------------------------------------------------------------------

\AtBeginDocument{%
  \begin{center}
    {\large\bfseries Fonds Alexandre Grothendieck --- cote n\textsuperscript{o}~\grothFolder}\\[2pt]
    {\itshape\grothFoldertitle}\\[4pt]
    {\small feuillets \grothFirst--\grothLast\ (lot \grothBatch)}\\[2pt]
    {\footnotesize\color{gray!60!black}Transcription établie d'après le fac-similé numérisé
     par l'Université de Montpellier.\\
     Les lectures incertaines sont soulignées, les illisibles notées [\,\ldots\,],
     les ajouts entre crochets.}
  \end{center}
  \vspace{1em}}

\endinput


\folder{150}
\batch{2}
\pages{21}{40}
\dating{1981-1982}
\watermark{Édition de démonstration}
\foldertitle{Espaces stratifiés et voisinages côniques (ou : déploiement des espaces stratifiés) : notes manuscrites (s.d.)}

\begin{document}

\section*{Germes de voisinages côniques, ou cylindriques. Transversalité}

\page{21}
\note{En haut à gauche, « 6 » de sa main : sixième feuillet de la suite « Espaces stratifiés (étude par voisinages coniques) » du lot 1 (feuillets 1) à 5), pages 11 à 19). Le titre est précédé du numéro de paragraphe (5) cerclé ; il reprend donc le paragraphe (5) de la page 20, biffée.}

(5)\note{numéro cerclé.} \emph{Germes de voisinages côniques, ou cylindriques.} \struck{\emph{Compatibilités}} \emph{Transversalité}.

Soit $Y\subset X$. Deux voisinages \struck{\ill{}} \add{$T,T'$} \uncertain{côniques} de $Y$ dans $X$ sont \emph{équivalents} s'il existe $\lambda,\lambda'\in\,]0,1[$ tels que $\lambda T=\lambda' T'$ (égalité \add{compatible} avec les structures côniques sur $Y$). C'est une relation d'équivalence. Une classe d'équivalence est un \emph{germe} de voisinage cônique. On dit qu'un $X_i\subset X$ est compatible avec un germe, représenté par $T$, s'il est \add{$\exists\lambda\in\,]0,1[$ tel qu'il soit} compatible avec $\lambda T$ (donc aussi avec $\lambda' T$ pour $0<\lambda'\le\lambda$) -- ce qui ne dépend pas du choix du représentant. Idem pour la notion de compatibilité stricte.

\marginal{$\varepsilon\in\,]0,1[$, donc par $\lambda T=\varepsilon\lambda' T'$}\note{note en travers dans la marge gauche, en face de « $\lambda T=\lambda'T'$ » ; lecture incertaine.}

Si $T$ et $T'$ sont équivalents, \struck{alors} $T$ est cylindrique ssi $T'$ l'est. On trouve ainsi la notion d'équivalence de voisinages cylindriques, et \add{celle} de germe de voisinages cylindriques. (Pour qu'il \uncertain{en} existe, il faut que $Y$ soit \uncertain{un sous-espace bordant}.)

\marginal{on dira aussi \ill{} $Y$ bordant $X$, \ill{} \ill{} \ill{} bordant \ldots}\note{note en travers dans la marge gauche, en face de la parenthèse ; lecture très incertaine.}

Item pour équisingularité d'un germe de voisinages côniques.

Soit $(\tilde X^{*},\dot T)$, où $\dot T\subset\tilde X^{*}$ muni d'un germe de voisinage cylindrique. Soit $\dot T\to Y$ application continue, alors \struck{\ill{}}
\[
X=\tilde X^{*}\amalg_{\dot T}Y\supset Y
\]
est muni d'un germe de voisinage \struck{cy} cônique, qui est équisingulier ssi $\dot T\to Y$ est fibrant.

\page{22}
\note{En haut à gauche, un chiffre surchargé, illisible ; la page continue le feuillet 6 (page 21).}

\marginal{NB \ill{} \ill{} \uncertain{Th.} \ill{} \ill{} \ill{} un \uncertain{voisinage} \ill{} \ill{} \ill{} \ill{} voisinage cylindrique \ill{} (catégories \ill{})}\note{Plusieurs lignes en travers dans l'angle supérieur gauche, serrées contre le texte ; on ne lit sûrement que « NB », « voisinage cylindrique », « catégories ».}

\ill{} \add{deux} catégories \struck{\ill{}} : \emph{a)} celle des \struck{couples} \add{systèmes} $(\tilde X^{*},\dot X,\mathcal E,Y,f)$, où $\tilde X^{*}$ esp. top., $\dot X$ partie fermée, $\mathcal E$ un germe de voisinages cylindriques de $\dot X$ dans $\tilde X^{*}$, $Y$ esp. top., $f:\dot X\to Y$ \uncertain{application} continue ; \emph{b)} celle des systèmes $(X,Y,\mathcal T)$, où $X$ esp. top., $Y$ partie fermée, $\mathcal T$ un germe de voisinages côniques de $Y$ dans $X$. \struck{Comme Alors} Alors le foncteur \add{ci-dessus} \uncertain{naturel} de \emph{a)} dans \emph{b)} est une équivalence de catégories.

Rappelons que
\[
X=\tilde X^{*}\amalg_{\dot X}Y .
\]

\marginal{NB $(\tilde X^{*},\dot X)$, germe de \ill{} \uncertain{cônique} de \ill{} $X$ \ill{} $(Y,\mathcal T)$}\note{note en travers dans la marge gauche ; lecture incertaine.}

\emph{NB} \struck{Une} \add{Une partie fermée $X'$} de $X$ revient donc à une partie $\tilde X'^{*}$ de $\tilde X^{*}$, $Y'$ de $Y$, ayant même image inverse dans $\dot X$, et $X'$ est fermée ssi $\tilde X'^{*}$ et $Y'$ le sont. \struck{Pour que $X'$ soit} \add{Pour que $X'$ soit} compatibles avec le germe de structure cônique $\mathcal T$, \struck{il faut et suffit qu'une description \ill{} que $X'$} il faut et suffit que l'on ait \emph{a)} $\tilde X'$ compatible avec le germe de structure cylindrique de $(\tilde X,\dot X)$. Il y a lieu alors de compléter $\tilde X'$ par $\overline{X'\cap X^{*}}$ ($=X^{*}=\tilde X\setminus\dot X$), et le germe de $\tilde X'_1$ au voisinage\note{Au-dessus de « $\tilde X'$ », en fin de ligne, il a écrit « $\tilde X'_1$ » ; la phrase continue sur la page suivante.}

\page{23}
\note{En haut à gauche, « 7 » de sa main.}

de $\dot X\subset\tilde X$ est déterminé par la connaissance de
\[
\tilde X'_1\cap\dot X=\dot X'_1 .
\]

\emph{Corollaire} Les parties \add{fermées} $X'$ de $X$ compatibles avec le germe de structure cônique, sont en corr. 1-1 avec les couples $(\tilde X'\subset\tilde X,\ Y'\subset Y)$ de parties \struck{fermées} fermées de $\tilde X$ et de $Y$, telles que
\begin{enumerate}
  \item[a)] $\tilde X'$ soit compatible avec le \add{germe de} structure cylindrique de $\tilde X,\dot X$ ;
  \item[b)] $f:\dot X\to Y$ applique \struck{\ill{}} $\dot X'\overset{\text{déf}}{=}\tilde X'\cap\dot X$ dans $Y'$.
\end{enumerate}
On aura
\[
X'=\underbrace{\mathrm{Im}\,(\tilde X'\ \text{par}\ \tilde X\to X)}_{X'\setminus Y\cap X'}\cup Y'
\]
\note{Sous l'accolade, « $Y\cap X'$ » est lui-même souligné d'une petite accolade marquée « $Y'$ ».}

Pour que $X'$ soit \struck{strictement compatible avec} \add{cônique strict}, \struck{le germe de structure cônique} il faut et suffit que $Y'=\mathrm{Im}\,(\dot X'\to Y)$.

Soit maintenant $X_0\subset X$ une partie bordante de $X$, et $\mathcal E$ une \add{germe de} structure cylindrique de $X_0$ dans $X$. On dit que $\mathcal E$ est compatible avec le germe de structure cônique $\mathcal T$ de $X$ autour de $Y$, si $\exists$ représentants $E$ \add{$\simeq X_0\times\mathbf{I}$} voisinage cylindrique de $X_0$, et $T$ voisinage cônique de $Y$ (avec structure cônique), tels que

a) $T\cap E$, $\dot T\cap E$, et $Y\cap E$ \struck{compatible} \add{soient} des parties cylindriques \struck{et $E\cap T$ soit stable par les deux types d'homothéties $x\mapsto\lambda.x$ que celles \ill{}}\ldots

\marginal{on dira aussi « $E$ est compatible » ou « \ill{} \ill{} » \ldots\ donc stable par $E$-homothéties par $\lambda\in[0,1]$}\note{notes en travers dans la marge gauche ; lecture incertaine.}

\note{Le bas de la page (depuis « et $E\cap T$ soit stable ») est barré de longues diagonales et repris en haut de la page 24. On y lit encore, sous les traits : « i.e. » $\dot T\cap E\simeq\dot T_0\times\mathbf{I}$, $T\cap E\simeq T_0\times\mathbf{I}$, $Y\cap E\simeq Y_0\times\mathbf{I}$, avec $T_0=X_0\cap T$, $Y_0=\dot X_0\cap Y$.}

\page{24}
\note{Sans numéro de sa main ; suite du feuillet 7 (page 23).}

\[
\begin{array}{c}
\dot T\cap E\simeq\dot T_0\times\mathbf{I}\\
\cap\\
T\cap E\simeq T_0\times\mathbf{I}\\
\cup\\
Y\cap E\simeq Y_0\times\mathbf{I}
\end{array}
\qquad\text{où}\qquad
\begin{array}{l}
\dot T_0=\dot T\cap X_0\\
T_0=T\cap X_0\\
Y_0=Y\cap X_0
\end{array}
\]
\note{Devant la deuxième ligne, un mot biffé (« \uncertain{i.e.} »).}

et si, symétriquement, \struck{$T\cap E$, $T\cap\dot E$, $X_0$, $T\cap E$}

b) $E\cap T$ $(=T\cap E)$, $\dot E\cap T$ \struck{$\simeq T_0\times\partial\mathbf{I}$}, $X_0\cap T$ $(=T\cap X_0=T_0)$ sont des \uncertain{ensembles} côniques de $T$, donc en particulier stables par $T$-homothéties. \struck{Le fait que} Notons que les trois ensembles précédents s'identifient, grâce à a), à : $T_0\times\mathbf{I}$, $T_0\times\lbrace1\rbrace$, $T_0\times\lbrace0\rbrace$. On suppose enfin

c) Les homothéties \struck{par} $\lambda_E$ et $\lambda_T$ \add{commutent} dans $T\cap E$, ce qui implique que l'application induite par $\lambda_T$ ($\lambda\in[0,1]$) sur $T_0\times\lbrace0\rbrace$ et $T_0\times\lbrace1\rbrace$ s'identifie, via \struck{homothétie} une opération de $T_0\to T_0$, \struck{\ill{}} soit $\lambda_{T_0}$, et que l'opération $\lambda_T$ sur $T\cap E\simeq T_0\times\mathbf{I}$ s'identifie à : \struck{\ill{}} $\lambda_{T_0}\times\mathrm{id}_{\mathbf{I}}$.

Donc les conditions a) b) c) équivalent à celle-ci :

\emph{Condition de compatibilité} $\exists$ homéomorphisme
\[
T\cap E\simeq T_0\times\mathbf{I}\qquad(\text{compatible avec } E\simeq X_0\times\mathbf{I})
\]
\struck{compatible avec structure} induisant
\[
\dot T\cap E\simeq\dot T_0\times\mathbf{I},\qquad Y\cap E\simeq Y_0\times\mathbf{I}
\]
tel que la structure cônique de $T$ sur $Y$ \add{induite} sur $T_E$ sur $Y_E$ \struck{soit} \add{soit} la structure \emph{cônique produit} \add{par $\mathbf{I}$,} d'une structure cônique de $T_0$ sur $Y_0\subset T_0$, \struck{par $\mathbf{I}$}

\page{25}
\note{En haut à gauche, « 8 » de sa main.}

\emph{NB} Cette condition implique que pour $\lambda_0,\lambda'_0\in[0,1]$, $\lambda_{0E}E$ et $\lambda'_{0T}T$ sont compatibles.

Si $T$ est aussi cylindrique, la condition \uncertain{envisagée} de compatibilité des deux structures cylindriques est symétrique.

\emph{Prop.} \add{(}Si $(X,Y,\mathcal T)$ est décrit par $(\tilde X,\dot X,\mathcal E_{\dot X},Y,\dot X\to Y)$\add{)}, alors les \add{couples} $(X_0,\mathcal E)$ \add{germes de} structures cylindriques compatibles avec le germe de structure cônique $\mathcal T$ correspondent aux \struck{couples} \add{systèmes}
\[
(\tilde X_0,\ \mathcal E_{\tilde X_0},\ Y_0,\ \mathcal E_{Y_0})
\]
où $\tilde X_0$ est une partie bordante de $\tilde X$, munie d'un germe de voisinages cylindriques $\mathcal E_{\tilde X_0}$, où $Y_0$ est une partie bordante de $Y$, munie d'un germe de \struck{voisin} voisinages cylindriques $\mathcal E_{Y_0}$, ces données satisfaisant aux conditions suivantes :

a) $\mathcal E_{\tilde X_0}$ et $\mathcal E_{\dot X}$ sont des \add{germes de} structures cylindriques \emph{compatibles} dans \add{$\tilde X$}, donc $\dot X_0\subset\dot X$ \add{($\dot X_0=\tilde X_0\cap\dot X$)} hérite d'un \struck{germe} germe de structure cylindrique (induit par $\mathcal E_{\tilde X_0}$) ;

b) Le morphisme $\dot X\xrightarrow{\ p\ }Y$ \struck{$(\dot X,\ldots)\to(Y,\ldots)$} est compatible avec $(\dot X_0,\mathcal E_{\dot X_0})$ et $(Y_0,\mathcal E_{Y_0})$, dans le sens \add{que $\dot X_0=p^{-1}(Y_0)$, et} qu'il existe des voisinages cylindriques respectifs de ces germes, \uncertain{munis de} \struck{\ill{}} $E_{\dot X}\simeq\dot X_0\times\mathbf{I}$, $E_Y\simeq Y_0\times\mathbf{I}$, tels que

\page{26}
\note{Sans numéro de sa main ; suite du feuillet 8.}

$p\,|\,\dot X_0\times\mathbf{I}$ induise $\dot X_0\times\mathbf{I}\to Y_0\times\mathbf{I}$, s'identifiant à $p_0\times\mathrm{id}_{\mathbf{I}}$, où $p_0:X_0\to Y_0$.\note{Il écrit ici $X_0$ sans point ; on attend $\dot X_0$.}

\struck{\emph{Compatibilité}} \emph{Transversalité} d'une famille finie $(B_i,E_i)$ de voisinages cylindriques de parties $B_i$ de $X$.

La condition s'exprime en disant que pour tout $J\subset I$, $I$ fini, la structure de $\bigcap_{i\in J}E_{B_i}\overset{\text{déf}}{=}E_J$, au-dessus de $B_J\overset{\text{déf}}{=}\bigcap_{i\in J}B_i$, \add{\uncertain{est isomorphe}} à celle de $B_J\times\mathbf{I}^J$, l'homothétie $\lambda_{E_j}$ \add{(pour $j\in J$)} de $E_J$ (ci-dessus) \uncertain{correspondant à}
\[
(b,(t_\alpha)_{\alpha\in J})\longmapsto(b,(t'_\alpha)),\qquad
t'_\alpha=t_\alpha\ \text{si}\ \alpha\neq j,\quad t'_\alpha=\lambda t_\alpha\ \text{si}\ \alpha=j .
\]

\marginal{et homéomorphisme $E_J\simeq B_J^{\mathbf{I}}$ \ill{} \ill{} \ldots\ $\dot B_J=\bigcap_{i\in J}\dot E_i$ \ldots}\note{notes en travers dans la marge gauche, en partie biffées ; lecture très incertaine.}

Cela peut aussi s'expliciter en exigeant que $E_J$ soit cylindrique dans chaque $E_{B_j}$ \ldots

Cette relation de transversalité \ill{} \uncertain{pour} un germe \ill{}

On dit que la famille précédente \emph{et} un voisinage cônique $T$ de $Y\subset X$ sont \emph{transversaux} si \add{\uncertain{de plus}} $\forall J$ \struck{\ill{}} comme dessus, \struck{compatible si la structure des $T\cap E_J$} :
\[
\begin{array}{lll}
E_J\cap T\simeq T_J\times\mathbf{I}^J & & (\text{où } T_J=T\cap B_J)\\
\qquad\cup & & \\
\dot E_J\cap\dot T\simeq\dot T_J\times\mathbf{I}^J & \text{induisant} & (\text{où } \dot T_J=\dot T\cap B_J)\\
E_J\cap Y\simeq Y_J\times\mathbf{I}^J & \text{induisant} & (\text{où } Y_J=Y\cap B_J)
\end{array}
\]
\note{En marge, en face de la première ligne : « \uncertain{induit par} $E_J\simeq B_J\times\mathbf{I}^J$ ». Dans la deuxième ligne il écrit $\dot E_J\cap\dot T$, dans la parenthèse finale de la troisième $Y\sim B_J$ pour $Y\cap B_J$.}

\page{27}
\note{En haut à gauche, « 9 » de sa main.}

avec $E_J\cap T$ cônique sur $T$, et la structure cônique \struck{induite} étant la structure produit par $\mathbf{I}^J$ d'une structure cônique sur $(T_J,Y_J)$, ayant $\dot T_J$ comme partie-bord des cônes sur $Y_J$\ldots

\section*{Relations avec stratifications de type $I$}

(6.)\note{numéro cerclé.} \emph{Relations avec stratifications de type $I$} ($I$ un ordonné fini).

Une telle stratification (\struck{\ill{}} \add{espace} $X$) est la donnée d'une famille croissante de parties fermées $X_i$. On dit que la stratification est \emph{précylindrique} si :

a) $\forall j'\leftrightarrow j$, $j$ successeur de $j'$, $X_{j'}\hookrightarrow X_j$ est \add{bordant} ;

b) \struck{$\forall j$} $\forall j$ fixé, $j'$ variable parmi les prédécesseurs, les \ill{} transversaux \ill{} dans $X_j$ \uncertain{$I_{:j}$}\,] \uncertain{La} structure de \ill{} un

c) $\forall i\leftrightarrow j$, \ill{} \ill{} \ill{} $k$ tel que $i\le k\le j$ est celle de l'un variété d'un \struck{simplexe} \add{(=$S_{ij}$)} -- ce qui signifie que si $S$ est l'ens. des prédécesseurs de $j$ dans $I_{i:j}$, i.e. des éléments minimaux de $I_{i:j}\setminus\lbrace j\rbrace$, et si pour tout $k\in I_{ij}$, on désigne par $s_k$ \struck{\ill{}} \add{la} partie $s_k\subset S$ formée \add{des} $J\in S$ tels que $J\ge s_k$ (donc $s_k=s_{kJ}$), l'application $k\longmapsto s_k$ (qui est décroissante) \struck{est un} est un anti-isomorphisme de $I_{ij}$ sur $\mathfrak P(S)$ \struck{\ill{}} \uncertain{ordonné} par inclusion.

\marginal{i.e. voisinages cylindriques \ill{} \ill{} transversaux \ldots\ $X_{j'}$ \ldots\ \emph{NB} Cela \ill{} de l'existence d'un simplexe \ldots\ On \ill{} \ill{} \ill{} \ill{} des \ill{} (c) dans \ill{}}\note{Longues notes en travers dans la marge gauche, en face de b) et c) ; lecture très incertaine. Dans la marge droite, un trait vertical ponctué de petits signes accompagne c). Tout ce passage est serré et surchargé : la lecture de b) et de la première ligne de c) est très incertaine.}

Combinant b) et c), on voit que la

\page{28}
\note{Sans numéro de sa main ; suite du feuillet 9.}

structure locale de $X_j$ le long d'un $X_i$ (avec la stratification induite de type $I_{ij}$) est « cubique » (coin de cube de dimension $n=\mathrm{card}\,S$).\note{À droite, un croquis : un coin de carré (deux axes), une courbe au-dessus et une demi-droite hachurée issue du sommet.}

\struck{La condition} c) \uncertain{implique} que si $j'$ est un prédécesseur de $j$, et si $k\in I$ est $\le j$ mais non $\le j'$, et s'il existe un $i$ majoré par $k$ et $j'$, \emph{alors} $k'=\inf(k,j')$ existe, et $k'$ est un prédécesseur de $k$.

\begin{tikzcd}
 & j' \arrow[r, leftrightarrow, "\text{préd}"'] \arrow[d, no head, "\geq" description] & j \arrow[d, no head, "\geq" description] \\
 & k' \arrow[r, leftrightarrow] & k \\
i \arrow[uur] \arrow[urr] & &
\end{tikzcd}
\note{Croquis de sa main : les traits verticaux portent le signe d'ordre (couché), les flèches horizontales ont une pointe à chaque bout, et deux traits partent de $i$ vers $j'$ et vers $k$.}

(Il suffit de voir que dans un ens. ordonné de la forme $\mathfrak P(F)$, les inf existent (l'intersection) et si $J'\subset J=J'\amalg\lbrace\alpha\rbrace$, et $k\in J$ mais non $k\subset J'$, alors \struck{\ill{}} $k=k'\amalg\lbrace\alpha\rbrace$ avec $k'=k\cap J'$, \ldots)

\marginal{$\sigma$ \uncertain{sur} $F$}\note{en marge gauche.}

Ceci posé, considérons le diagramme

\begin{tikzcd}
X_{j'} \arrow[r] & X_j \\
X_{k'} \arrow[u] \arrow[r] & X_k \arrow[u]
\end{tikzcd}

où les flèches verticales, par l'hyp. \ill{}, sont des plongements bordants. Ceci posé, on voit \add{(par b, c)} que \struck{$(X_k\to X)$} $(X_k,X_{k'})\to(X_j,X_{j'})$

\section*{Stratifications et voisinages côniques}

\page{29}
\note{Sans numéro de sa main.}

7. \emph{Stratifications} et \emph{\struck{\ill{}} voisinages côniques} : le \uncertain{jeu} \add{des \uncertain{récurrences}}\ldots

$(X_i)_{i\in I}$, famille croissante de parties de $X$, satisfaisant
\[
(14)\qquad
\begin{array}{ll}
a) & \forall i,j\in I,\ \text{on a}\ X_i\cap X_j=\bigcup_{\alpha\le i,j}X_\alpha .\\[4pt]
b) & X=\bigcup_{i\in I}X_i
\end{array}
\]
\note{Numéro (14) tel quel : il a déjà servi à la page 18 (lot 1) pour une autre formule.}

\marginal{Cette condition \uncertain{généralise} \ldots\ $\bigcap X_i=\bigcup_{k\in I,\ k\le\ldots}X_k$ \ldots\ \ill{} \ill{} \ill{} \ldots}\note{notes en travers dans la marge gauche, en face de (14) ; lecture très incertaine.}

\emph{Transversalité} : un syst. de parties bordantes $B_\alpha$ \add{$B_\alpha$} de $X$, muni de germes de cylindres \uncertain{associés} (mutuellement transversaux entre eux) : si on \ill{} \ill{} par des cylindres \ill{} \uncertain{représentant} les germes que \ill{} \ill{} les $X_i$ soient transversaux $E_{B_\alpha}$, \ill{} \ill{} \uncertain{supposons} de l'$E_{B_\alpha}$ -- cela implique que l'isom.
\[
E_J\cap X_i\simeq(X_{iJ}\times\mathbf{I}^J)\qquad\text{(induit par } E_J\simeq B_J\times\mathbf{I}^J\text{, où } X_{iJ}=X_i\cap B_J\text{)}.
\]
\struck{\ill{}}

\marginal{$J$ \ill{} \ill{} pour \ill{} d'indices}\note{en marge gauche.}

\struck{\ill{}} Soit $I_0\subset I$ une partie de $I$ formée d'éléments minimaux, \struck{\ill{}} et soit
\[
X_{I_0}=\bigcup_{i\in I_0}X_i=\coprod_{i\in I_0}X_i .
\]

\struck{\ill{}} La donnée d'un germe \add{strict} \struck{de structure} \add{de} $X_{I_0}$ équivaut : d'un voisinage cônique équivaut à celle d'un syst. de voisinages côniques $T_i$ ($i\in I_0$) mutuellement disjoints des $X_i$. -- celle d'un germe de voisinage cônique de $X$, à celle d'une famille de germes de vois. côniques des $X_i$ dans $X$.\note{« strict » est ajouté au-dessus de « germe » ; les deux « équivaut » sont tels quels, la phrase a été reprise sans être remise en ordre. « de $X$ » : on attend $X_{I_0}$.}

\page{30}
\note{En haut à gauche, « 10 » de sa main. La page reprend la suite de la page 28 (fin de c) de (6.)).}

\struck{Soit} « transversal » i.e. $\exists$ voisinage cylindrique \add{de $X_{j'}$,} \struck{et un} de $X_{k'}$ dans $X_k$, tel que l'application $X_k\to X_j$ induise un \add{morphisme} \struck{homéomorphisme} sur les cylindres (compatible avec les structures cylindriques) et de plus
\[
E_{k;k}=p_{jk}^{-1}(E_{j';j}) .
\]
(avec $X_{j'}$ dans les $X_j$)\note{Indices lus tels quels : on attend $E_{k';k}$.}

\emph{Système de voisinages cylindriques} \struck{\ill{}} \emph{compatible avec une stratification} : tel que dans b), les $E_{j';j}$ ($j$ fixé, $j'$ prédécesseur) soient compatibles. Notion correspondante pour les germes de voisinages cylindriques.

\emph{Stratification cylindrique} : stratification munie d'une telle « structure cylindrique », en termes de germes de voisinages cylindriques.

\struck{\emph{Mettre}}

\emph{NB} Une stratification précylindrique \add{\uncertain{resp.} cylindrique} \struck{est sûrement} \uncertain{est} « équisingulière » \ill{} un sens \ill{} plus fort que \ill{} \ill{} (C) \uncertain{suite}.

Mais j'ai oublié une condition, que j'exigerai pour les stratifications :
\[
\forall i,j\in I,\ \text{on a}\ X_i\cap X_j=\bigcup_{k\le i,j}X_k .
\]
Corollaire : si $X_i\cap X_j\neq\emptyset$, alors $i,j$ ont un minorant $k$ (tel que $X_k\neq\emptyset$\ldots)

\marginal{\emph{NB} si $i,j$ ont un \ill{} $h$, \uncertain{$\inf(i,j)$} \ill{} $X_i\cap X_j=X_h$ \ill{} \ill{} \ill{} minorant}\note{note en travers dans la marge gauche ; lecture incertaine.}

\page{31}
\note{En haut à gauche, « 11 » de sa main ; suite de la page 29.}

Prenons $T_{I_0}=\coprod_{\beta\in I_0}T_\beta$ voisinage \struck{\ill{}} \add{cônique} strict de $X_{I_0}$, supposons le strict transversal aux $(B_i,\mathcal E_i)$, et \struck{aussi} que les $X_i$ ($i\in I^{*}\overset{\text{déf}}{=}I\setminus I_0$) soient str. admissibles pour $T_{I_0}$.

Notons que en vertu de (14), si $\beta\in I_0$ donc minimal dans $I$, on a pour $i\in I$\ldots
\[
(14')\qquad X_i\cap X_\beta=
\begin{cases}
\emptyset & \text{si } i\not\ge\beta\\
X_\beta & \text{si } i\ge\beta\ \text{(cas trivial)}
\end{cases}
\]
Soit $I_\beta$ \add{$\subset I^{*}$} l'ens. des $i\in I$ tels que $i>\beta$. On a donc
\[
(15)\qquad i\in I^{*}\setminus I_\beta\Longrightarrow X_i\cap T_\beta=\emptyset
\]
(car $X_i\cap X_\beta=\emptyset$ et $X_i$ admissible pour $T_\beta$),
\[
(16)\qquad i\in I_\beta\Longrightarrow X_i\supset X_\beta
\]
Donc on a
\[
(16')\qquad X_i\cap T_\beta=\mathbf{I}.\underbrace{X_i\cap\dot T_\beta}_{\dot X_i}\cup X_\beta
\]

Ainsi, le système
\[
(17)\qquad\bigl(X,\ (X_i)_{i\in I},\ (B_\alpha,\mathcal E_\alpha)_{\alpha\in R},\ (T_\beta)_{\beta\in I_0}\bigr)
\]
se reconstitue à l'aide de \ldots\note{Numéros (15), (16) tels quels : le lot 1 a déjà une formule (15), page 20 (biffée). Le reste de la page -- une première écriture du système (18) : $X^{*}=X\setminus\overset{\circ}{T}_{I_0}$, $(\dot X_\beta)$ avec $\dot X_\beta=\dot T_\beta$, $(B^{*}_\alpha,\mathcal E^{*}_\alpha)$ avec $B^{*}_\alpha=B_\alpha\cap X^{*}$, $\mathcal E^{*}_\alpha=\mathcal E_\alpha|X^{*}$, $(X_{\beta\alpha},\mathcal E_{\beta\alpha})_{\beta\in I_0}=(B^{*}_\alpha,\mathcal E^{*}_\alpha)|X_\beta$, $(\dot X_\beta\to X_\beta)_{\beta\in I_0}$ -- est barré de longues diagonales et repris à la page 32.}

\page{32}
\note{En haut à gauche, « 12 » de sa main.}

\[
(18)\quad\left\lbrace
\begin{array}{l}
\overbrace{X^{*},\ (\dot X_\beta)_{\beta\in I_0},\ (B^{*}_\alpha,\mathcal E^{*}_\alpha)_{\alpha\in\Lambda}}^{\text{syst. de cylindres-bords transversaux}},\ (X^{*}_i)_{i\in I^{*}}\ ;\ (X_\beta)_{\beta\in I},\ (B_{\beta,\alpha},\mathcal E_{\beta,\alpha})_{\alpha\in\Lambda,\ \beta\in I_0},\\[6pt]
(\dot X_\beta\xrightarrow{\ p_\beta\ }X_\beta)_{\beta\in I_0}
\end{array}\right\rbrace
\]
avec $X^{*}=X\setminus\overset{\circ}{T}_{I_0}$, $\dot X_\beta=\dot T_\beta$, $(B^{*}_\alpha,\mathcal E^{*}_\alpha)=(B_\alpha,\mathcal E_\alpha)|X^{*}$, $X^{*}_i=X_i\cap X^{*}$, $(B_{\beta,\alpha},\mathcal E_{\beta,\alpha})=(B_\alpha,\mathcal E_\alpha)|X_\beta$, de \uncertain{façon} suivante :\note{Ces égalités sont écrites sous les termes de (18), reliées par des « = » verticaux. L'indice $\beta\in I$ de $(X_\beta)$ est tel quel ; on attend $\beta\in I_0$. L'accolade supérieure est de sa main.}

\[
(19)\quad\left\lbrace
\begin{array}{l}
\hat X=X^{*}\amalg_{\dot X=\coprod\dot X_\beta}\underbrace{\mathcal C\bigl(p=\textstyle\coprod p_\beta\bigr)}_{\coprod_\beta\mathcal C(p_\beta)},\\[10pt]
T_\beta=\mathcal C(p_\beta)\\[4pt]
X_\beta\ \text{pour}\ \beta\in I_0\ \text{clair}\\[4pt]
X_i\ \text{pour}\ i\in I^{*}\ \text{est donné par}\\[4pt]
\qquad\bar X_i=X^{*}_i\cup\bigcup_{\beta\in I_0,\ \beta\le i}\bigl(\mathbf{I}.\dot X_{i\beta}\cup X_\beta\bigr),\qquad \dot X_{i\beta}\overset{\text{déf}}{=}X^{*}_i\cap\dot X_\beta\\[10pt]
B_\alpha=B^{*}_\alpha\cup\bigcup_{\beta\in I_0}\mathbf{I}.\dot B_\alpha,\qquad \dot B_\alpha\overset{\text{déf}}{=}B^{*}_\alpha\cap\dot X_\beta\\[10pt]
\mathcal E_\alpha\ \text{défini par}\ \mathcal E^{*}_\alpha\ \text{et les}\ (B_{\beta,\alpha},\mathcal E_{\beta\alpha})
\end{array}\right.
\]
\note{Dans la marge gauche de (19) : « les $\dot X_\beta$ disjoints ». Sur la ligne de $X_i$, une première écriture « $X_i=\mathbf{I}.\dot X_i$ \ldots » est biffée, et dans la formule de $\bar X_i$ un terme en $\dot X_{\beta}$ est biffé d'un gribouillis ; l'indice « $\beta\le i$ » remplace une première écriture biffée. Le $\hat X$ (lettre coiffée) et le $\bar X_i$ sont tels quels.}

\emph{Conditions sur (18)}
\[
(20)\quad
\begin{array}{ll}
(a) & \text{Les } X_\beta \text{ disjoints}\\
(b) & \text{Chaque } \dot X_\beta \text{ transverse : } (B^{*}_\alpha,\mathcal E^{*}_\alpha)\ \Longrightarrow\ \text{le syst. induit dans } (\dot X_\beta,\ldots)\ (\dot B^{\alpha},\dot{\mathcal E}^{\alpha})\\
(c) & (\dot X_\beta,\underbrace{\dot B_{\beta\alpha},\dot{\mathcal E}_{\beta\alpha}}_{\text{induit sur }\dot X_\beta\text{ par }(B_\alpha,\mathcal E_\alpha)})\longrightarrow(X_\beta,B_{\beta,\alpha},\mathcal E_{\beta,\alpha})\ \text{est transverse}\\
(d) & \text{Les } X^{*}_i \text{ transverses aux } \bigl((\dot X_\beta,\mathcal E_\beta)_{\beta\in I_0},(B^{*}_\alpha,\mathcal E^{*}_\alpha)_{\alpha\in\Lambda}\bigr)
\end{array}
\]
\note{Les lettres (a) à (d) sont cerclées. Une accolade réunit (a) et (b), d'où part la flèche vers « le syst. induit dans $\dot X_\beta$ » (lecture incertaine, en bord de page), et une autre, en marge gauche, réunit (b) à (d) ; le numéro (20) est porté en marge. En face de (c), à droite : « \uncertain{est transverse} ».}

\page{33}
\note{Sans numéro de sa main ; suite du feuillet 12.}

Il faut de plus des conditions qui assurent, à partir des données (18), que $i\mapsto X_i$ est croissant, et que la condition (14)
\[
X_i\cap X_j=\bigcup_{k\le i,j}X_k .
\]
Pour \struck{les} exprimer ces conditions, introduisons pour tout $\beta\in I_0$ l'ens.
\[
(21)\qquad I^{*}_\beta=\lbrace i\in I^{*}\mid i>\beta\rbrace
\]
donc la donnée de $(I,I')$ \add{(I ens. ordonné)} équivaut à celle de $I^{*}$, $I_0$ \add{(ens. ordonnés)} et de la famille des $I^{*}_\beta\subset I^{*}$ ($\beta\in I_0$) \add{(\ill{} $j\in I^{*}$)}.\note{« $(I,I')$ » tel quel ; on attend $(I,I_0)$.}

Pour \struck{exprimer} \add{exprimer} que $i\le j$ implique $X_i\subset X_j$, on distingue le cas où $i\in I_0$ \add{($=\beta$)}, \struck{\ill{}} où c'est clair sur (19), et celui où $i\in I^{*}$, \struck{où} où c'est clair aussi sur (19), compte tenu de la condition suivante évidente :
\[
(e)\qquad i,j\in I^{*},\ i\le j\Longrightarrow X^{*}_i\subset X^{*}_j
\]
\struck{Pour on doit exiger} On doit avoir aussi, manifestement
\[
(f)\qquad i,j\in I^{*}\Longrightarrow X^{*}_i\cap X^{*}_j=\bigcup_{h\le i,j}X^{*}_h
\]

Pour la relation
\[
X_i\cap X_j=\bigcup_{h\le i,j}X_h\qquad(i,j\in I)
\]
il est clair qu'il suffit de prouver le cas où $i,j$ ne sont pas comparables. Quand un des deux, soit $\beta$, est $\in I_0$, alors cette relation \struck{est} signifie que pour $\beta\in I_0$, $i\in I^{*}\setminus I_\beta$, on a $X_\beta\cap X_i=\emptyset$, ce qui résulte de (19).

\page{34}
\note{En haut à gauche, « 13 » de sa main.}

Si $i,j\in I^{*}$, il suffit \add{(grâce à (f))} de vérifier que pour tout $\beta$, on a
\[
X_i\cap X_j\cap T_\beta=\bigcup_{h\le i,j}X_h\cap T_\beta
\]
(il suffit de prouver $\subset$, car $\supset$ trivial). Si $X_i\cap X_j\cap T_\beta=\emptyset$ c'est trivial, on peut donc supposer $X_i\cap X_j\cap T_\beta\neq\emptyset$ \struck{i.e. $X_i\cap X_j\cap$} donc $X_i\cap T_\beta\neq\emptyset$, $X_j\cap T_\beta\neq\emptyset$ i.e. $i,j\in I_\beta$. Alors \struck{cela} cela résulte de (f) induit sur $\dot T_\beta=\dot X_\beta\subset X^{*}$, et de la formule (19). Donc

\emph{Th.} Les systèmes (17), satisfaisant aux conditions \struck{(14)} (14) des stratifications, et aux conditions de transversalité explicitées plus haut, équivalent aux systèmes (18), satisfaisant aux conditions (20), en plus des conditions (e) (f) exprimant que les $(X^{*}_i)_{i\in I^{*}}$ forment une stratification de $X^{*}$.

\section*{Structure cônique sur un espace stratifié}

8. \emph{Structure cônique sur un espace stratifié}

Pour $X$ stratifié de type \struck{$I$} $I$ par des $(X_i)_{i\in I}$, alors des sous-espaces bordants donnés $(B_\alpha)_{\alpha\in\Lambda}$ et des \struck{germes} germes de \add{voisinages} cylindriques \uncertain{sur} $B_\alpha$, tels que \struck{\ill{} mutuellement} \add{mutuellement} transverses, et les $(B_\alpha)$ soient transverses aux $X_i$, on va \struck{construire} définir

\page{35}
\note{Sans numéro de sa main ; suite du feuillet 13 (page 34).}

par récurrence sur $\mathrm{card}\,I$ la notion de donnée \struck{de} cônique sur \struck{cette base}
\[
(X,(X_i)_{i\in I},(B_\alpha,\mathcal E_\alpha)_{\alpha\in\Lambda}).
\] \struck{On posera} Si $I=\emptyset$, on pose la structure vide (il y en a une et une seule). Si $I\neq\emptyset$, on prend $I_0$ ens. des éléments minimaux de $I$, et une donnée cônique sur par définition la donnée

(a)\note{lettre cerclée.} d'un \struck{\ill{}} \add{vois. con. strict} $T_{I_0}$ de \struck{\ill{}} $X_{I_0}=\bigcup_{\beta\in I_0}X_\beta$, transversal aux $(B_\alpha,\mathcal E_\alpha)$ et \struck{\ill{} admissible} côniquement transversal aux $X_i$ ($i\in I$).

-- ce qui permet donc d'introduire la structure (19), et en particulier
\[
(22)\qquad \underbrace{X^{*},\ (X^{*}_i)_{i\in I^{*}=I\setminus I_0}}_{\text{stratification de } I^{*}},\ 
\underbrace{\bigl((B_\alpha,\mathcal E_\alpha)_{\alpha\in\Lambda},\ (\dot X_\beta,\mathcal E_\beta)_{\beta\in I_0}\bigr)}_{\text{famille transversale}}
\]
\note{Sous la seconde accolade, en cinq lignes : « famille transversale de sous-espaces bordants \uncertain{« cylindrisés »} de $X$, transversaux aux $X_i$ ». Devant $B_\alpha$, une lettre biffée (un point surmonté, peut-être $\dot B_\alpha$ ou $\mathcal E$).}

en plus de
\[
(23)\qquad \bigl((X_\beta)_{\beta\in I_0},\ (B_{\beta\alpha},\mathcal E_{\beta\alpha})_{\beta\in I_0,\ \alpha\in\Lambda}\bigr)
\]

On doit avoir de plus :

(b)\note{lettre cerclée.} Une structure cônique sur (22) (ce qui a un sens, par récurrence\ldots)

\page{36}
\note{En haut à gauche, « 14 » de sa main.}

\struck{\ill{}} \emph{Remarques} 1) Avec des hypothèses de « non singularité » sur la donnée initiale $(X,(X_i)_{i\in I},(B_\alpha,\mathcal E_\alpha)_{\alpha\in\Lambda})$, il doit être vrai qu'il existe une structure cônique. Si la stratification \struck{\ill{}} est équisingulière le long de tout
\[
X_i^{\circ}\overset{\text{déf}}{=}X_i\setminus\bigcup_{j<i}X_j ,
\]
alors il doit y avoir une structure cônique équisingulière (avec la définition ad-hoc évidente) -- il suffit de savoir construire un tube équisingulier autour de $X_{I_0}$\ldots\note{Sous l'indice $j<i$, une seconde écriture « $j\le i$ ».}

2) A priori, si on a une suite \add{str.} croissante de parties $\emptyset\subset I_0\subset I_1\subset\ldots\subset I_{n-1}\subset I_n=I$, avec \struck{$I_i$ formé d'él. minimaux} \add{$I_i\setminus I_{i-1}$} formé d'él. minimaux de $I\setminus I_{i-1}$, alors on pourrait définir \add{par récurrence} une notion de structure cônique pour cette « filtration » de $I$. Mais on peut montrer que cette notion ne dépend pas, en fait, essentiellement, de la filtration choisie.

3) Au lieu d'une famille de $(B_\alpha,\mathcal E_\alpha)_{\alpha\in\Lambda}$, on peut supposer donnée sur $X$ une stratification cylindrique \struck{de type} (\ill{} \ill{} \ill{} \ill{} celui où l'ensemble ordonné qui indexe la stratification cylindrique est $\mathfrak P(\Lambda)^{\circ}$.

\page{37}
\note{Sans numéro de sa main ; suite du feuillet 14.}

Il n'y a pas de difficulté essentielle au traitement fait ici. Plus tard, on généralisera cette notion de stratification cylindrique (globale) en une notion de stratification cylindrique locale -- et \ill{} \ill{} \ill{}\ldots

\section*{Description d'une stratification cônique de type $I$}

9. \emph{Description d'une \struck{structure} \add{stratification} cônique de type $I$} ($I$ un ordonné fini).

Pour simplifier, on prend \add{d'abord} le cas où on ne se donne pas de stratifications cylindriques sur $X$ \uncertain{par dessus le marché}.

\struck{\ill{}}\note{Un mot biffé ; le reste de la page est blanc.}

\page{38}
\note{En haut à gauche, « 15 » de sa main.}

\emph{Stratification \struck{équisingulière} \add{équisingulière} de type $I$} ($I$ un ens. ordonné fini) : \emph{présentation standard}. On suppose donné une « \struck{catégorie} catégorie topologique modèle » $\mathcal M$ (\struck{\ill{}} par ex. pour \uncertain{fixer les} idées, pour que ce puisse être aussi celle des espaces topologiques\ldots). La présentation est donnée par un foncteur
\[
\Sigma:\mathrm{Drap}(I)\longrightarrow\mathcal M,\qquad (i_0,\ldots,i_n)\longmapsto\Sigma_{i_0\ldots i_n}
\]
avec les propriétés suivantes

(1) Les \struck{\ill{}} morphismes
\[
\Sigma_{i_0\ldots i_n}\longrightarrow\Sigma_{i_0,\ldots,i_p}
\]
(relatifs : une inclusion de drapeaux $d\hookrightarrow d'$ telle que $d$ soit un segment initial de $d'$) sont \emph{fibrants}${}^{*}$ \struck{(2)} (il suffit de le supposer pour $\Sigma_{i_0\ldots i_n}\to\Sigma_{i_0\ldots i_{n-1}}$)

\marginal{${}^{*}$ C'est trop faible, si on veut l'équisingularité -- il faut supposer que pour $d$ fixé, et pour $d'=d\amalg\delta$, $\delta$ variable (i.e. $\delta$ est donné comme drapeau quelconque des $I^{\,\ill{}}$, \ill{} \ill{} \ill{}) la famille des $\Sigma_{d'}$ au-dessus de $\Sigma_d$ est loc. triviale}\note{Note de renvoi, écrite en travers dans la marge droite, rattachée par l'astérisque à « fibrants » ; lecture incertaine dans le détail.}

(2) Les morphismes
\[
\Sigma_{i_0\ldots i_n}\longrightarrow\Sigma_{i_0\ldots\widehat{i_p}\ldots i_n},\qquad (p\neq n)
\]
\struck{$\Sigma_{d'}\to\Sigma_d$} sont des \emph{plongements-bord} (NB si $p=n$, c'est un \add{morphisme} fibrant) \add{i.e. plongements \uncertain{bords} cylindriques}.\note{« $p\neq n$ » est entouré.}

\marginal{NB cette condition \ill{} \ill{} \ill{} \ill{} \uncertain{équisingulière} \ldots\ On veut de plus \ill{} une structure cylindrique \ill{} \ill{} \ill{} conditions (2), \ill{} \ill{} \ill{}}\note{notes en travers dans la marge gauche ; lecture très incertaine.}

\emph{Cor.} Soit $d\hookrightarrow d'$ tel que $d$ et $d'$ aient même \ill{} \ill{} \ill{} \ill{} \ill{}, alors $\Sigma_{d'}\to\Sigma_d$ est un plongement.

\page{39}
\note{Sans numéro de sa main ; suite du feuillet 15. La page est très chargée : notes en haut, à droite et, en travers, dans toute la marge gauche.}

\marginal{${}^{*}$ Il suffit \uncertain{déjà} que l'on ait que si $d\cap d'$ soit segment initial de $d'$ et contienne le plus grand élément de $d$, ce qui signifie par (D), $\Sigma_d\to\Sigma_\delta$ \ill{} et $\Sigma_{d''}\to\Sigma_d$ est fibrant}\note{Note en tête de page, séparée du texte par un trait, rattachée par l'astérisque à « sans plus » ; lecture incertaine. (D) est cerclé.}

(3) Soient $d,d'$ deux drapeaux tels que \add{sans plus}${}^{*}$ le dernier élément $i$ de $d$ soit le premier élément de $d'$ et considérons $d''=d\cup d'$ et le diagramme

\begin{tikzcd}
\Sigma_{d''} \arrow[r, hook] \arrow[d] & \Sigma_{d'} \arrow[d] \\
\Sigma_d \arrow[r, hook] & \Sigma_i
\end{tikzcd}

Ce dernier est \emph{cartésien}. (Pour $d$ ou $d'$ \ill{} i.e. réduit à $\lbrace i\rbrace$, cette condition est triviale)

\marginal{Il suffit de le supposer quand $d$ et $d'$ sont de longueur 2 (\ill{} \ill{} \ill{}). On en déduit plus généralement que lorsque $d,d'$ sont deux drapeaux et que $d''=d\cup d'$ soit un drapeau, \ill{} $d''=d\amalg_\delta d'$ \ill{} \ill{}, alors (D) $\Sigma_{d''}\to\Sigma_{d'}$, $\Sigma_d\to\Sigma_\delta$ \ill{} cartésien.}\note{Note encadrée à droite, en face du diagramme, écrite serrée ; (D) est cerclé et accompagne un petit carré $\Sigma_{d''}\to\Sigma_{d'}$, $\Sigma_d\to\Sigma_\delta$ ; lecture incertaine.}

De plus, il faut supposer que lorsque $d$ est de longueur 2, alors $\Sigma_{d'}\to\Sigma_\delta$ est \emph{transverse} aux sous-espaces cylindriques $\Sigma_d$ de $\Sigma_{\ill{}}$ -- et \ill{} pour ces structures cylindriques \ill{} \ill{}\note{Passage serré, ajouté en face de la fin de (3) ; lecture incertaine.}

Cela revient donc à \ill{} \ill{} $\Sigma_{\ill{}}$ \ill{} \ill{} construire les $\Sigma_d$, \uncertain{précisément} \ill{} \uncertain{connaissant} les $\Sigma_i$, $\Sigma_{ij}$ \add{($i<j$)} et les morphismes

\begin{tikzcd}
\Sigma_{ij} \arrow[r, hook] \arrow[d] & \Sigma_j \\
\Sigma_i &
\end{tikzcd}

\struck{Pour} (NB, la verticale est fibrante, et l'horizontale un morphisme-bord). Pour avoir les morphismes de transition entre les $\Sigma_{d'}$ (construits par récurrence sur long $d$, \uncertain{pour} long 0 et 1 étant déjà donné) il faut se donner les morphismes
\[
\underbrace{\Sigma_{ij}\times_{\Sigma_j}\Sigma_{jk}}_{\overset{\text{déf}}{=}\Sigma_{ijk}}\longrightarrow\Sigma_{i,k}
\]
satisfaisant aux conditions\ldots

\marginal{\ill{} $\mathrm{Dr}(I)$ \ldots\ a) $\Sigma_d\cap\Sigma_{d'}=\emptyset$ si \ill{} \ldots\ \uncertain{$\Sigma_{i_0\ldots i_n}=\Sigma_{i_0}\cap\ldots$} \ldots\ b) $\Sigma_{i_0\ldots i_n}$ \ill{} \ldots\ ($\Leftarrow$ $\Sigma_d\cap\Sigma_{d'}=\Sigma_{d\cup d'}$) \ldots\ c) On \ill{} \ill{} cylindriques \ill{} \ill{} \ill{} transverses \ill{}}\note{Longues notes en travers dans toute la marge gauche, sur une douzaine de lignes serrées ; seuls quelques fragments se lisent, dans une lecture très incertaine.}

\page{40}
\note{En haut à gauche, « 16 » de sa main.}

La donnée des $\Sigma_i$ équivaut : à la donnée d'un $\Sigma_{*}\in\mathcal M$ ($\Sigma_{*}=\coprod\Sigma_i$) et d'une partition de type $I$ de $\Sigma_{*}$ i.e. d'un morphisme $\Sigma_{*}\to I_{\mathcal M}$. Celle des $\Sigma_{ij}$ équivaut : celle d'un objet $\Sigma_{**}$ \struck{\ill{}} au-dessus de $I\times I$ \struck{$\mathrm{Drap}_2(I)$}. \struck{$I\times I$} (on pose $\Sigma_{ij}=\emptyset$ si $i\not\le j$, $\Sigma_{ii}=\Sigma_i$)

La donnée des

\begin{tikzcd}
\Sigma_{ij} \arrow[r, hook] \arrow[d] & \Sigma_j \\
\Sigma_i &
\end{tikzcd}

équivaut : celle de deux morphismes « source » et « but »

\begin{tikzcd}
\Sigma_{**} \arrow[r, "\underline b"] \arrow[d, "\underline s"'] & \Sigma_{*} \\
\Sigma_{*} &
\end{tikzcd}

compatibles avec les $\mathrm{pr}_1$ et $\mathrm{pr}_2$ de $I\times I$ dans $I$ (\struck{\emph{NB} la \ill{}} Comme la bigraduation \add{de type $I$} de $\Sigma_{**}$ est déduite de \struck{celle} la graduation de type $I$ de $\Sigma_{*}$, et de $\underline b$ et $\underline s$) -- où $\underline b$ est \add{loc.\textsuperscript{t}} un morphisme bord (i.e. est une immersion \struck{bord}) et $\underline s$ est une fibration. On a de plus
\[
\underline i:\Sigma_{*}\longrightarrow\Sigma_{**}
\]
« morphisme identité », dont le composé avec $\underline b$ et $\underline s$ est l'identité. La donnée des $\Sigma_{ijk}\to\Sigma_{ik}$ revient :

\end{document}
